\section{Pretraining data：规模之外还有来源、顺序与外部记忆}

上一章解释了模型如何计算，本章追问梯度究竟来自什么分布。对训练数据分布 $q(x)$，模型实际优化
\[
\mathbb{E}_{x\sim q}\big[\ell_{\theta}(x)\big].
\]
改变 $q$ 就改变模型反复看到的事实、语言、风格与推理模式。Token 数只是一个维度；质量、覆盖、重复率、顺序、检索接口和参数增长 schedule 都会改变有效监督。课程用四个研究项目把“数据很重要”拆成可检验的问题。

\subsection{数据不是油桶，而是训练目标的一部分}

把数据叫作“新石油”容易让人只想到数量。更准确的说法是：数据定义经验风险的采样分布。若网页语料中某种模板、错误事实或语言占比过高，优化器会把它们当作稳定规律；若高价值样本很少，即使模型容量足够，梯度也未必能反复强化该能力。因此 data curation、mixture 与 schedule 是 architecture 的外部组成部分。

\lecturefigure{slide-035.jpg}{数据对 AI 的影响：质量、规模、算力与能力共同耦合}{CS25 V6 official deck, p.~35}

\lecturefigure{slide-036.jpg}{数据策略项目一：个体儿童输入与双语学习}{CS25 V6 official deck, p.~36}

\lecturefigure{slide-037.jpg}{数据策略项目二：RAG 计算分配与 curriculum-guided scaling}{CS25 V6 official deck, p.~37}

\teachervoice{00:16:20--00:20:23，老师把数据研究分成四类干预：改变来源、改变语言混合、把知识放到外部检索、以及让数据难度与模型深度共同变化。讲义提醒：这些项目回答不同因果问题，不能合成一句“更好数据更好模型”。}

\begin{knowledgebox}{读图：四个项目改变的是不同变量}
Baby Scale 改变\textbf{家庭/儿童数据源}；Bilingual BabyLM 改变\textbf{语言与说话人混合结构}；RAG scaling 改变\textbf{参数记忆与外部记忆的 compute allocation}；CGLS 改变\textbf{模型深度与样本难度的时间顺序}。先确定 intervention，再看 metric，最后才能讨论可迁移结论。
\end{knowledgebox}

\begin{importantbox}{数据策略的五个旋钮}
\begin{enumerate}
  \item \textbf{Selection}：哪些样本进入 corpus；
  \item \textbf{Mixture}：语言、领域和任务如何配比；
  \item \textbf{Order}：简单/困难、旧/新数据何时出现；
  \item \textbf{Externalization}：知识放进参数还是 retrieval store；
  \item \textbf{Co-scaling}：数据 schedule 是否与参数、层数和 compute 同步变化。
\end{enumerate}
这些旋钮都会改变梯度，不能被视为训练前的一次性清洗。
\end{importantbox}

这四类项目也展示了数据研究的基本因果设计。若改变数据来源，应尽量固定 tokenizer、参数量、训练 token 和优化器；若改变语言 mixture，应同时报告每种语言的测试分布；若改变 retrieval budget，要把额外 context 与索引成本计入；若改变 layer schedule，则需控制总参数更新量和训练 compute。否则所谓“data effect”可能只是模型更大、训练更久或评估更接近训练分布。课程没有把这些实验写成一条万能定律，反而提供了一个方法论：每次只明确一组 intervention，并把尚未控制的 confound 写进结论边界。

\subsection{本章小结}

Pretraining data 决定的不只是模型“看过什么”，还包括样本被看到的频率、顺序、上下文和可检索性。后面四个案例将分别回答：小而自然的数据能学到什么，双语是否互相干扰，RAG 应分走多少预算，以及模型成长能否与 curriculum 同步。

